\section{A scaled formulation}\label{sec:setup}

The proof separates the thermodynamic limit from the small gas parameter.
We first work in units in which each unit cell contains many particles,
while the exclusion distance is much smaller than their typical separation.
The physical change of units will be made in Section~\ref{sec:physical}.

Let $K\ge100$ be an integer and let $\Lambda_K=(\R/K\Z)^3$.
Write
\begin{equation}\label{eq:scaled-parameters}
 V=K^3,\qquad D=N/V,\qquad r=\alpha/D,\qquad
 \alpha\in[\alpha_0/2,2\alpha_0],\qquad 0<T\le D^{-1},\qquad \beta=T^{-1}.
\end{equation}
The number $\alpha_0\ge2$ will be fixed sufficiently large; subsequently
$D$ will be taken sufficiently large. For $j\ge1$, let
\[
 \Omega_j=\{X\in\Lambda_K^j:d(x_i,x_k)>r\text{ for }i\ne k\},
 \qquad q_j[f]=\sum_{i=1}^j\int_{\Omega_j}|\nabla_i f|^2.
\]
Here $d$ is Euclidean torus distance. The form domain is
$H^1_0(\Omega_j)$, restricted to symmetric functions for the bosonic
Hamiltonian $H_j$. Functions are extended by zero to $\Lambda_K^j$.
This convention is used for partial traces, multipliers, and integral
kernels throughout the paper.

For a bounded measurable function $W\ge0$, set
\[
 H_j^{(-W)}=H_j-\sum_{i=1}^jW(x_i),\qquad
 Z_j^{(-W)}=\Tr e^{-\beta H_j^{(-W)}},\qquad
 F_j^{(-W)}=-T\log Z_j^{(-W)}.
\]
The trace is on the symmetric space. Omit the superscript when $W=0$,
and put $Z_0=1$, $F_0=0$. For a symmetric trace-one state $\sigma$ of
finite kinetic energy, write
\[
 S(\sigma)=-\Tr\sigma\log\sigma,\qquad
 \mathcal F_j^{(-W)}(\sigma)
 =\Tr H_j^{(-W)}\sigma-TS(\sigma).
\]
Energy expectations mean form expectations. The entropy bound in
Lemma~\ref{lem:entropy} implies that these entropies are finite.
Gibbs' variational principle gives
$\mathcal F_j^{(-W)}(\sigma)\ge F_j^{(-W)}$.
The insertion construction in Section~\ref{sec:energy} also verifies that
the relevant hard-core form domains, through particle number $2N$, are
nonzero for the parameters used here.

The unit cells are
$B_v=v+[-1/2,1/2)^3$, periodically interpreted, with
$v\in\mathcal L_K=(\Z/K\Z)^3$. The function $n_v(X)$ counts positions
in $B_v$. For a set of cells $A\subset\mathcal L_K$, distances to $A$
mean distances to its centers. Changing between the union of cells and
their centers changes the estimates below only by absolute constants.
The lattice distance $d_\infty$ is the periodic maximum distance on
$\mathcal L_K$. Unless stated otherwise, $c,C>0$ denote constants
independent of $K,D,T$ in the ranges under consideration; they may change
from line to line and depend on parameters already fixed. Dependence on
the fixed interval for $\alpha$ is allowed after $\alpha_0$ is chosen.

\subsection{Positive kernels and retained particles}

Let $H_j^{\rm lab}$ be the labelled Dirichlet operator on $\Omega_j$,
let $K_s^{(j)}$ be its heat kernel, and let $P_j$ be the orthogonal
projection onto symmetric functions. When $j=N$, denote the kernel of
$e^{-sH_N^{\rm lab}}P_N$ by $B_s^{\rm sym}$. This operator is the
bosonic semigroup extended by zero on the orthogonal complement of the
symmetric subspace.
Explicitly it is the average of $K_s^{(N)}$ over permutations of one
endpoint. It is nonnegative and obeys the semigroup law with integration
over ordered configurations, without any additional factorial.

The labelled kernel is the free Brownian kernel, with generator
$\Delta$, restricted to paths whose pair distances stay strictly above
$r$. One way to establish this exact boundary convention is to exhaust
$\Omega_j$ by relatively compact smooth open sets. Their Dirichlet forms
increase in domain to $H^1_0(\Omega_j)$, giving strong semigroup
convergence. A continuous path that stays in $\Omega_j$ on a compact time
interval has compact image there and is captured by the exhaustion.
Monotone convergence of killed path measures proves the representation.
The same reasoning, or bounded-potential Feynman--Kac, treats the fields
above. This is the open-domain argument used in \citet[Section~4]{HS}.

Heat traces are finite: the zero-extension form embedding and the
min--max principle compare the Dirichlet eigenvalues to those of the free
torus operator. Kernel identities are needed only almost everywhere.
Splitting a semigroup at positive times justifies diagonal trace formulas
without selecting values of a general kernel on a null set.

Partial traces preserve the remaining hard-core conditions and the kinetic
expectations of the retained coordinates. The form-domain justification is
given with the entropy estimates in Section~\ref{sec:energy}, before these
marginals are used to compare particle numbers.

\subsection{The estimate to be proved}

Let $\Gamma_N=Z_N^{-1}e^{-\beta H_N}$, let
$\gamma_N^{(1)}=N\Tr_{2,\ldots,N}\Gamma_N$, and let $u_0=V^{-1/2}$.
The occupation fraction is
\[
 B_{\rm occ}=N^{-1}\langle u_0,\gamma_N^{(1)}u_0\rangle.
\]

\begin{theorem}[Scaled occupation bound]\label{thm:scaled}
There are constants $\alpha_0\ge2$, $D_*<\infty$, and $c_*>0$ such that
$B_{\rm occ}\ge c_*$ whenever the parameters satisfy
\eqref{eq:scaled-parameters}, $D\ge D_*$, and $K$ is sufficiently large.
The lower threshold for $K$ can be chosen locally uniformly for $D$ in
$[D_*,\infty)$, uniformly in $\alpha\in[\alpha_0/2,2\alpha_0]$ and
$0<T\le D^{-1}$.
\end{theorem}

Sections~\ref{sec:energy} and \ref{sec:spatial} establish free-energy
estimates and joint bounds on the position counts. Section~\ref{sec:slab}
extends those bounds to a short interval of particle paths. The remaining
sections compare configurations by moving selected labels along paths of
usable cells. This yields Theorem~\ref{thm:scaled} in
Section~\ref{sec:comparison}; the final change of units proves
Theorem~\ref{thm:main}.
